% -*- coding: utf-8 -*-
% !TEX program = xelatex
\documentclass[plain,noanswer,medmath]{../../template/jnuexam}
\usepackage{../../template/kysx}

\begin{document}


\examtitle{2016}{1}


\makepart{选择题}{1～8小题，每小题4分，共32分}


\begin{problem}
若反常积分$\int_0^{+\infty} \frac{1}{x^a(1+x)^b}\dx$收敛，则\pickout{}
\begin{abcd}
\item $a < 1$且$b > 1$．
\item $a > 1$且$b > 1$．
\item $a < 1$且$a + b > 1$．
\item $a > 1$且$a + b > 1$．
\end{abcd}
\end{problem}


\begin{problem}
已知函数$f(x) = \begin{cases}
  2(x - 1), & x < 1, \\
     \ln x, & x \ge 1,
\end{cases}$，则$f(x)$的一个原函数是\pickout{}
\begin{abcd}
\item $F(x) = \begin{cases}
                 (x - 1)^2, & x < 1 \\
  x\left(\ln x - 1 \right), & x \ge 1
\end{cases}$．
\item $F(x) = \begin{cases}
                     (x - 1)^2, & x < 1 \\
  x\left(\ln x + 1 \right) - 1, & x \ge 1
\end{cases}$．
\item $F(x) = \begin{cases}
                     (x - 1)^2, & x < 1 \\
  x\left(\ln x + 1 \right) + 1, & x \ge 1
\end{cases}$．
\item $F(x) = \begin{cases}
                     (x - 1)^2, & x < 1 \\
  x\left(\ln x - 1 \right) + 1, & x \ge 1
\end{cases}$．
\end{abcd}
\end{problem}


\begin{problem}
若$y_1 = \left(1 + x^2 \right)^2 - \sqrt{1 + x^2}$，$y_2 = \left(1 + x^2 \right)^2 + \sqrt{1 + x^2}$
是微分方程$y' + p(x)y = q(x)$的两个解，则$q(x) =$\pickout{}
\begin{abcd}
\item $3x\left(1 + x^2 \right)$．
\item $- 3x\left(1 + x^2 \right)$．
\item $\frac{x}{1 + x^2}$．
\item $- \frac{x}{1 + x^2}$．
\end{abcd}
\end{problem}


\begin{problem}
已知函数$f(x) = \begin{cases}
            x, & x \le 0, \\
  \frac{1}{n}, & \frac{1}{n + 1} < x \le \frac{1}{n},n = 1,2, \cdots ,
\end{cases}$则\pickout{}
\begin{abcd}
\item $x = 0$是$f(x)$的第一类间断点．
\item $x = 0$是$f(x)$的第二类间断点．
\item $f(x)$在$x = 0$处连续但不可导．
\item $f(x)$在$x = 0$处可导．
\end{abcd}
\end{problem}


\begin{problem}
设$A$，$B$是可逆矩阵，且$A$与$B$相似，则下列结论错误的是\pickout{}
\begin{abcd}
\item $A^T$与$B^T$相似．
\item $A^{-1}$与$B^{-1}$相似．
\item $A + A^T$与$B + B^T$相似．
\item $A + A^{-1}$与$B + B^{-1}$相似．
\end{abcd}
\end{problem}


\begin{problem}
设二次型$f\left(x_1,x_2,x_3 \right) = x_1^2 + x_2^2 + x_3^2 + 4x_1x_2 + 4x_1x_3 + 4x_2x_3$，
则$f\left(x_1,x_2,x_3 \right) = 2$在空间直角坐标下表示的二次曲面为\pickout{}
\begin{abcd}
\item 单叶双曲面．
\item 双叶双曲面．
\item 椭球面．
\item 柱面．
\end{abcd}
\end{problem}


\begin{problem}
设随机变量$X\sim N\left(\mu ,\sigma^2 \right)$（$\sigma > 0$），
记$p = P\left\{X \le \mu + \sigma^2 \right\}$，则\pickout{}
\begin{abcd}
\item $p$随着$\mu$的增加而增加．
\item $p$随着$\sigma$的增加而增加．
\item $p$随着$\mu$的增加而减少．
\item $p$随着$\sigma$的增加而减少．
\end{abcd}
\end{problem}


\begin{problem}
随机试验$E$有三种两两不相容的结果$A_1,A_2,A_3$，且三种结果发生的概率均为$\frac{1}{3}$，
将试验$E$独立重复做$2$次，$X$表示$2$次试验中结果$A_1$发生的次数，
$Y$表示$2$次试验中结果$A_2$发生的次数，则$X$与$Y$的相关系数为\pickout{}
\begin{abcd}
\item $- \frac{1}{2}$．
\item $- \frac{1}{3}$．
\item $\frac{1}{3}$．
\item $\frac{1}{2}$．
\end{abcd}
\end{problem}


\makepart{填空题}{9～14小题，每小题4分，共24分}


\begin{problem}
$\lim_{x \to 0} \frac{\int_0^x t\ln \left(1 + t\sin t \right)\dt}{1 - \cos x^2} =$ \fillin{}．
\end{problem}


\begin{problem}
向量场$A\left(x,y,z \right) = \left(x + y + z \right)i + xyj + zk$的旋度
$\operatorname{rot} A =$ \fillin{}．
\end{problem}


\begin{problem}
设函数$f(u,v)$可微，$z = z(x,y)$由方程$(x + 1)z - y^2 = x^2 f(x - z,y)$确定，
则$\dz\big|_{(0,1)}=$ \fillin{}．
\end{problem}


\begin{problem}
设函数$f(x) = \arctan x - \frac{x}{1 + ax^2}$，且$f'''(0) = 1$，则$a =$ \fillin{}．
\end{problem}


\begin{problem}
行列式$\begin{vmatrix}
  \lambda &      -1 &       0 & 0 \\
        0 & \lambda &      -1 & 0 \\
        0 &       0 & \lambda & -1 \\
        4 &       3 &       2 & \lambda + 1
\end{vmatrix} =$ \fillin{}．
\end{problem}


\begin{problem}
设$x_1,x_2, \cdots ,x_n$为来自总体$N\left(\mu ,\sigma^2 \right)$的简单随机样本，
样本均值$\overline{x} = 9.5$，参数$\mu$的置信度为$0.95$的双侧置信区间的置信上限为$10.8$，
则$\mu$的置信度为$0.95$的双侧置信区间为 \fillin{}．
\end{problem}


\makepart{解答题}{15～23小题，共94分}


\begin{problem}[points=10]
已知平面区域$D = \left\{(r,\theta)\middle| 2 \le r \le 2(1+\cos\theta),
- \frac{\pi}{2} \le \theta \le \frac{\pi}{2} \right\}$，\goodbreak 计算二重积分$\iint_D x\dx\dy$．
\end{problem}


\begin{problem}[points=10]
设函数$y(x)$满足方程$y'' + 2y' + ky = 0$，其中$0 < k < 1$．
\begin{enumerate}
\item 证明：反常积分$\int_0^{+\infty} y(x)\dx$收敛；
\item 若$y(0) = 1$，$y'(0) = 1$，求$\int_0^{+\infty} y(x)\dx$的值．
\end{enumerate}
\end{problem}


\begin{problem}[points=10]
设函数$f(x,y)$满足$\frac{\pd f(x,y)}{\pdx} = (2x + 1)\e^{2x - y}$，
且$f(0,y) = y + 1$．$L_t$是从点$(0,0)$到点$(1,t)$的光滑曲线，
计算曲线积分$I(t) = \int_{L_t} \frac{\pd f(x,y)}{\pdx}\dx + \frac{\pd f(x,y)}{\pdy}\dy$，
并求$I(t)$的最小值．
\end{problem}


\begin{problem}[points=10]
设有界区域$\Omega$由平面$2x + y + 2z = 2$与三个坐标平面围成，$\Sigma$为$\Omega$整个表面的外侧，
计算曲面积分$I = \iint_{\Sigma}\left(x^2 + 1 \right)\dy\dz - 2y\dz\dx + 3z\dx\dy$．
\end{problem}


\begin{problem}[points=10]
已知函数$f(x)$可导，且$f(0) = 1$，$0 < f'(x) < \frac{1}{2}$，
设数列$\left\{x_n \right\}$满足$x_{n + 1} = f(x_n)$（$n = 1,2, \cdots$），证明：\par
\begin{enumerate*}
\item 级数$\mathop{\sum}\limits_{n = 1}^{\infty}(x_{n + 1} - x_n)$绝对收敛；
\item $\lim_{n \to \infty} x_n$存在，且$0 < \lim_{n \to \infty} x_n < 2$．
\end{enumerate*}
\end{problem}


\begin{problem}[points=11]
设矩阵$A = \begin{pmatrix}
   1 & -1 & -1 \\
   2 &  a & 1 \\
  -1 &  1 & a
\end{pmatrix}$，$B = \begin{pmatrix}
       2 & 2 \\
       1 & a \\
  -a - 1 & -2
\end{pmatrix}$．当$a$为何值时，方程$AX = B$无解、有唯一解、有无穷多解？在有解时，求解此方程．
\end{problem}


\begin{problem}[points=11]
已知矩阵$A = \begin{pmatrix}
  0 & -1 & 1 \\
  2 & -3 & 0 \\
  0 &  0 & 0
\end{pmatrix}$．
\begin{enumerate}
\item 求$A^{99}$；
\item 设$3$阶矩阵$B = (\alpha ,\alpha_2,\alpha_3)$满足$B^2 = BA$，
      记$B^{100} = (\beta_1,\beta_2,\beta_3)$，
      将$\beta_1,\beta_2,\beta_3$分别表示为$\alpha_1,\alpha_2,\alpha_3$的线性组合．
\end{enumerate}
\end{problem}


\begin{problem}[points=11]
设二维随机变量$(X,Y)$在区域
$D = \left\{(x,y) \middle| 0 < x < 1, x^2 < y < \sqrt{x} \right\}$
上服从均匀分布，令$U = \begin{cases}
  1, & X \le Y; \\
  0, & X > Y.
\end{cases}$
\begin{enumerate}
\item 写出$(X,Y)$的概率密度；
\item 问$U$与$X$是否相互独立？并说明理由；
\item 求$Z = U + X$的分布函数$F(z)$．
\end{enumerate}
\end{problem}


\begin{problem}[points=11]
设总体$X$的概率密度为$f(x,\theta) = \begin{cases}
  \frac{3x^2}{\theta^3}, & 0 < x < \theta, \\
  0,                     & \text{其他},
\end{cases}$其中$\theta \in (0,+\infty)$为未知参数，
${X_1},{X_2},{X_3}$为来自总体$X$的简单随机样本，
令$T = \max\{X_1,X_2,X_3\}$．\par
\begin{enumerate*}
\item 求$T$的概率密度；
\item 确定$a$，使得$aT$为$\theta $的无偏估计．
\end{enumerate*}
\end{problem}


\end{document}
